\documentclass[10pt,a4paper]{amsart}
\usepackage{amsmath,amssymb,amsthm,graphicx}
% Retained preamble for the Stacks Project chapter {CH} (shared project preamble).
% The chapter file pinned-original-source.tex only carries \input{{preamble}} on its first line,
% so no preamble text lives inside the chapter itself.  The Stacks Project uses one single shared
% preamble for every chapter: the snapshot stages it three times (sources/stacks-chow/preamble.tex,
% sources/stacks-etale-cohomology/preamble.tex and focused-seeds/AGP/sources/stacks-snapshot-20261011/
% preamble.tex) and those copies are identical in size.  The block below is copied mechanically,
% line for line, from sources/stacks-chow/preamble.tex lines 14-18 and 154-217: the Xy-pic setup and
% the native theorem declarations together with the mathematical macros.  Nothing was added, removed
% or redefined; the wrapper class and the amsmath/amssymb/amsthm/graphicx base packages are supplied
% by the assembler.
% For commutative diagrams we use Xy-pic
% Chapter: stacks-cohomology.tex -- Cohomology of Sheaves.
\usepackage[all]{xy}

% We use 2cell for 2-commutative diagrams.
\xyoption{2cell}

% Theorem environments.
%
\theoremstyle{plain}
\newtheorem{theorem}[subsection]{Theorem}
\newtheorem{proposition}[subsection]{Proposition}
\newtheorem{lemma}[subsection]{Lemma}

\theoremstyle{definition}
\newtheorem{definition}[subsection]{Definition}
\newtheorem{example}[subsection]{Example}
\newtheorem{exercise}[subsection]{Exercise}
\newtheorem{situation}[subsection]{Situation}

\theoremstyle{remark}
\newtheorem{remark}[subsection]{Remark}
\newtheorem{remarks}[subsection]{Remarks}

\numberwithin{equation}{subsection}

% Macros
%
\def\lim{\mathop{\mathrm{lim}}\nolimits}
\def\colim{\mathop{\mathrm{colim}}\nolimits}
\def\Spec{\mathop{\mathrm{Spec}}}
\def\Hom{\mathop{\mathrm{Hom}}\nolimits}
\def\Ext{\mathop{\mathrm{Ext}}\nolimits}
\def\SheafHom{\mathop{\mathcal{H}\!\mathit{om}}\nolimits}
\def\SheafExt{\mathop{\mathcal{E}\!\mathit{xt}}\nolimits}
\def\Sch{\mathit{Sch}}
\def\Mor{\mathop{\mathrm{Mor}}\nolimits}
\def\Ob{\mathop{\mathrm{Ob}}\nolimits}
\def\Sh{\mathop{\mathit{Sh}}\nolimits}
\def\NL{\mathop{N\!L}\nolimits}
\def\CH{\mathop{\mathrm{CH}}\nolimits}
\def\proetale{{pro\text{-}\acute{e}tale}}
\def\etale{{\acute{e}tale}}
\def\QCoh{\mathit{QCoh}}
\def\Ker{\mathop{\mathrm{Ker}}}
\def\Im{\mathop{\mathrm{Im}}}
\def\Coker{\mathop{\mathrm{Coker}}}
\def\Coim{\mathop{\mathrm{Coim}}}

% Boxtimes
%
\DeclareMathSymbol{\boxtimes}{\mathbin}{AMSa}{"02}

%
% Macros for moduli stacks/spaces
%
\def\QCohstack{\mathcal{QC}\!\mathit{oh}}
\def\Cohstack{\mathcal{C}\!\mathit{oh}}
\def\Spacesstack{\mathcal{S}\!\mathit{paces}}
\def\Quotfunctor{\mathrm{Quot}}
\def\Hilbfunctor{\mathrm{Hilb}}
\def\Curvesstack{\mathcal{C}\!\mathit{urves}}
\def\Polarizedstack{\mathcal{P}\!\mathit{olarized}}
\def\Complexesstack{\mathcal{C}\!\mathit{omplexes}}
% \Pic is the operator that assigns to X its picard group, usage \Pic(X)
% \Picardstack_{X/B} denotes the Picard stack of X over B
% \Picardfunctor_{X/B} denotes the Picard functor of X over B
\def\Pic{\mathop{\mathrm{Pic}}\nolimits}
\def\Picardstack{\mathcal{P}\!\mathit{ic}}
\def\Picardfunctor{\mathrm{Pic}}
\def\Deformationcategory{\mathcal{D}\!\mathit{ef}}
\usepackage{cleveref}
\crefname{originallemma}{original lemma}{original lemmas}
\Crefname{originallemma}{Original lemma}{Original lemmas}
\makeatletter
\newlabel{homology-lemma-homotopy-complex-complexes}{{0FKH}{source}{}{sourcepointer-1}{}}
\newlabel{homology-lemma-homotopy-complex-complexes@cref}{{[originallemma][0][]0FKH}{[1][0][]source}{}{}{}}
\makeatother
\crefname{originallemma}{original lemma}{original lemmas}
\Crefname{originallemma}{Original lemma}{Original lemmas}
\makeatletter
\newlabel{lemma-godement-resolution}{{0FKS}{source}{}{sourcepointer-2}{}}
\newlabel{lemma-godement-resolution@cref}{{[originallemma][0][]0FKS}{[1][0][]source}{}{}{}}
\makeatother
\title{Godement resolution of a bounded below complex by a flasque complex, stalkwise a homotopy equivalence}
\author{The Stacks Project Authors}
\date{stacks-project:cohomology.tex:local-snapshot-20261011; GNU FDL 1.2 or later}
\begin{document}
\maketitle
\noindent\textit{Editorial note.} All delivered mathematical text is the author's own, copied from the pinned Stacks Project chapter file cohomology.tex; no statement or proof step was rewritten, reordered or removed.\par\medskip
\noindent\textit{Editorial note.} The author's device of resolving complexes of modules by treating them as objects of an abelian category is written out at the start of the proof and the surrounding construction is complete as written.\par\medskip
\noindent\textit{Editorial note.} The counted spans carry 2 navigation stubs and no bibliography citation.\par\medskip
\noindent\textit{Editorial note.} The preamble is the chapter's retained minimal preamble (sources/stacks-cohomology/retained-preamble.tex lines 1-81), a verbatim excerpt of the native preamble.\par\medskip
\noindent\textit{Editorial note.} The label lemma-godement-resolution-bounded-below is defined only in cohomology.tex, so no cross-chapter restatement had to be weighed for this card.\par\medskip
\noindent\textit{Allowed standard prerequisites (editorial).}\begin{itemize}
\item Godement resolution of a sheaf of modules by flasque modules. Same chapter, Lemma lemma-godement-resolution, whose argument is reproduced for complexes at proof lines 7098-7100. Supplies the resolution pattern 0 to F to f\_*f\textasciicircum{}*F to f\_*f\textasciicircum{}*f\_*f\textasciicircum{}*F to ... that is being lifted to complexes.
\item Homotopy equivalence of complexes of complexes passes to the total complex. Homology, Lemma homology-lemma-homotopy-complex-complexes, invoked at proof lines 7136-7137. Supplies the stalkwise homotopy equivalence of the totalized complex G.
\item Pushforward from a discrete space preserves flasqueness. Same chapter, used at proof lines 7048-7051 in the sibling lemma and again implicitly here. Supplies that each term of the resolution is a flasque module.
\end{itemize}\medskip
\ifdefined\hypertarget\hypertarget{sourcepointer-1}{}\fi
\noindent\textit{Source reference.} Original lemma 0FKH: sources/stacks-chow/tags.txt, source lines 19462--19462. Pointer outside the counted span, invoked once for context; the load-bearing argument of the counted statement is the author's proof reproduced in full in the proof span.\par\medskip
\ifdefined\hypertarget\hypertarget{sourcepointer-2}{}\fi
\noindent\textit{Source reference.} Original lemma 0FKS: sources/stacks-cohomology/pinned-original-source.tex, source lines 7015--7015. Pointer outside the counted span, invoked once for context; the load-bearing argument of the counted statement is the author's proof reproduced in full in the proof span.\par\medskip
\begin{lemma}
\label{lemma-godement-resolution-bounded-below}
Let $(X, \mathcal{O}_X)$ be a ringed space. Let
$\mathcal{F}^\bullet$ be a bounded below complex of
$\mathcal{O}_X$-modules. There exists a quasi-isomorphism
$\mathcal{F}^\bullet \to \mathcal{G}^\bullet$
where $\mathcal{G}^\bullet$ be a bounded below complex of flasque
$\mathcal{O}_X$-modules and for all $x \in X$ the
map $\mathcal{F}^\bullet_x \to \mathcal{G}^\bullet_x$
is a homotopy equivalence in the category of complexes
of $\mathcal{O}_{X, x}$-modules.
\end{lemma}

\begin{proof}
Let $\mathcal{A}$ be the category of complexes of $\mathcal{O}_X$-modules
and let $\mathcal{B}$ be the category of complexes of $\mathcal{O}_X$-modules.
Then we can apply the discussion above to the adjoint functors
$f^*$ and $f_*$ between $\mathcal{A}$ and $\mathcal{B}$.
Arguing exactly as in the proof of
Lemma \ref{lemma-godement-resolution}
we get a resolution
$$
0 \to
\mathcal{F}^\bullet \to
f_*f^*\mathcal{F}^\bullet \to
f_*f^*f_*f^*\mathcal{F}^\bullet \to
f_*f^*f_*f^*f_*f^*\mathcal{F}^\bullet \to \ldots
$$
in the abelian category $\mathcal{A}$ such that each term of each
$f_*f^*\ldots f_*f^*\mathcal{F}^\bullet$ is a flasque
$\mathcal{O}_X$-module and such that for all $x \in X$ the
map
$$
\mathcal{F}^\bullet_x[0] \to \Big(
(f_*f^*\mathcal{F}^\bullet)_x \to
(f_*f^*f_*f^*\mathcal{F}^\bullet)_x \to
(f_*f^*f_*f^*f_*f^*\mathcal{F}^\bullet)_x \to \ldots
\Big)
$$
is a homotopy equivalence in the category of complexes of complexes
of $\mathcal{O}_{X, x}$-modules. Since a complex of complexes is the
same thing as a double complex, we can consider the induced map
$$
\mathcal{F}^\bullet \to
\mathcal{G}^\bullet =
\text{Tot}(
f_*f^*\mathcal{F}^\bullet \to
f_*f^*f_*f^*\mathcal{F}^\bullet \to
f_*f^*f_*f^*f_*f^*\mathcal{F}^\bullet \to \ldots
)
$$
Since the complex $\mathcal{F}^\bullet$ is bounded below, the
same is true for $\mathcal{G}^\bullet$ and in fact each term
of $\mathcal{G}^\bullet$ is a finite direct sum of
terms of the complexes $f_*f^*\ldots f_*f^*\mathcal{F}^\bullet$
and hence is flasque. The final assertion of the lemma
now follows from
Homology, Lemma \ref{homology-lemma-homotopy-complex-complexes}.
Since this in particular shows that
$\mathcal{F}^\bullet \to \mathcal{G}^\bullet$
is a quasi-isomorphism, the proof is complete.
\end{proof}

\end{document}
